\subsection{Por que dimensionar um sorteio no catálogo}
O catálogo já fornece as distribuições dos metadados. O sorteio teve outra finalidade prática: limitar o número de consultas para recuperar letras e, ao mesmo tempo, preservar uma seleção de origem conhecida para medir a cobertura. Como gênero explica uma parcela importante da variação de vários atributos, os 82 rótulos foram usados como estratos. Adotou-se alocação aproximadamente proporcional para manter a composição de gênero e evitar concentrar o desenho em apenas uma variável.

Sejam $N=1.159.764$ o tamanho do catálogo, $N_h$ o total do gênero $h$, $W_h=N_h/N$ e $n$ o tamanho a determinar. Para cada candidato a $n$, calculou-se $nW_h$, tomou-se sua parte inteira e distribuíram-se as vagas restantes pelos maiores restos. A busca exigiu ao menos duas observações por estrato. Após fixar a alocação $n_h$, foi realizado um sorteio aleatório simples sem reposição em cada gênero, independentemente dos demais, com semente 20261106. A probabilidade de inclusão é $\pi_i=n_h/N_h$ no estrato $h$ e o peso amostral é $w_i=1/\pi_i$ \cite{cochran1977,lohr2022,statcan2010}.

\subsection{Critérios e cálculo do tamanho}
Para uma variável $z$, a média ponderada por estrato e sua variância sob o desenho são
\begin{equation}
\widehat{\bar z}_{\mathrm{est}}=\sum_{h=1}^{82}W_h\bar z_h,\qquad
V_d(\widehat{\bar z}_{\mathrm{est}})=
\sum_{h=1}^{82}W_h^2\left(\frac{1}{n_h}-\frac{1}{N_h}\right)S_h^2,
\label{eq:desenho}
\end{equation}
onde $\bar z_h$ é a média no sorteio e $S_h^2$ é a variância finita do gênero, com divisor $N_h-1$. A expressão incorpora a correção de população finita \cite{cochran1977,lohr2022}. Para um indicador binário com proporção populacional $P_h$ no estrato, $S_h^2=N_h P_h(1-P_h)/(N_h-1)\leq N_h/[4(N_h-1)]$. O limite usa o cenário de maior variância, de aproximadamente metade dos registros em cada categoria do indicador.

Foram adotadas duas metas de planejamento. Uma proporção global, considerada individualmente, deveria ter margem nominal de até 1 ponto percentual a 95\% pela aproximação normal. Para as doze médias, de popularidade, ano e dez atributos de áudio e duração, definiu-se margem de até 0,05 desvio-padrão do respectivo campo, com nível conjunto nominal de 95\%. A escala relativa torna o critério comparável entre variáveis em unidades diferentes. O ajuste de Bonferroni usa a cauda $0,05/(2\cdot12)$ em cada margem. Com $V_{\max}$ calculada pelo limite binário, $S_j$ o desvio-padrão do campo $j$ no catálogo e $V_{d,j}$ sua variância de desenho, selecionou-se o menor tamanho examinado que atendesse simultaneamente a
\begin{equation}
z_{0{,}975}\sqrt{V_{\max}}\leq 0{,}01,\qquad
\max_{1\le j\le 12}\frac{z_{1-0{,}05/(2\cdot12)}\sqrt{V_{d,j}}}{S_j}\leq 0{,}05,
\end{equation}
em que $z_q$ é o quantil de ordem $q$ da normal padrão. Os momentos dos campos vieram da descrição do catálogo. A busca sequencial é apropriada porque as alocações inteiras por maiores restos mudam em degraus à medida que $n$ aumenta.

\begin{table}[H]\centering\small
\caption{Sensibilidade do dimensionamento à meta para uma proporção global.}\label{tab:dimensionamento}
\setlength{\tabcolsep}{5pt}
\begin{tabular}{@{}lrrr@{}}\toprule
Meta (p.p.) & Tamanho & Limite obtido (p.p.) & Maior margem média (DP) \\\midrule
2,00 & 3.195 & 1,73148 & 0,04999 \\
1,50 & 4.254 & 1,49984 & 0,04331 \\
1,00 & 9.526 & 0,99997 & 0,02887 \\\bottomrule
\end{tabular}
\nota{A meta para a família de doze médias é 0,05 DP em todos os cenários.}
\end{table}

A Tabela~\ref{tab:dimensionamento} mostra que, ao adotar a meta de 1 ponto percentual, o primeiro tamanho que satisfez ambos os critérios foi \textbf{9.526 faixas}. A margem para proporção foi a condição mais restritiva nesse cenário. A alocação final variou de 5 a 179 faixas por gênero. As metas se aplicam à descrição do sorteio antes da vinculação. O desenho não quantifica a precisão por gênero, quantas letras serão encontradas ou o poder para detectar ganhos preditivos.

\subsection{Verificação descritiva do sorteio}
A média ponderada da popularidade no sorteio foi de 18,40, ante 18,38 no catálogo. A mediana sem ponderação foi de 15 em ambos. Para as doze médias, o maior desvio padronizado da média ponderada em relação ao catálogo foi 0,0190. A Figura~\ref{fig:amostra} mostra a proximidade das médias e a comparação entre as dispersões observadas. Na Figura~\ref{fig:amostraperiodos}, a maior diferença entre proporções por período é de 0,405 ponto percentual. São verificações das marginais observadas, incluindo variáveis que não definiram os estratos.

\figura[.39\textheight]{comparacao_medias_dispersoes}{Médias e dispersões no catálogo e no sorteio inicial.}{fig:amostra}{As diferenças de médias são expressas em desvios-padrão do catálogo.}
\figura[.22\textheight]{comparacao_periodos}{Composição por período no catálogo e no sorteio inicial.}{fig:amostraperiodos}{Períodos definidos pelo ano informado no catálogo.}

\subsection{Vinculação das faixas ao LRCLIB}
O cruzamento teve como entrada as faixas Spotify selecionadas e como fonte de candidatos o dump do LRCLIB descrito na Seção~\ref{sec:fontes}. Como os IDs pertencem a sistemas diferentes, a busca usa créditos de artista, título e duração. O índice de metadados facilita a busca sem carregar letras de todos os candidatos. As chaves derivadas uniformizam Unicode, caixa, espaços e sinais tipográficos usuais. Preservaram-se os campos de origem e os qualificadores informativos de versão. A análise de participações explícitas busca evitar que uma diferença de crédito torne invisível um candidato compatível.

O processo verificou título e artista compatíveis e diferença de até dez segundos quando as duas durações eram válidas, após converter milissegundos para segundos. Possíveis homônimos e versões tiveram tratamento específico. Depois da seleção por metadados, letras sincronizadas foram interpretadas nos formatos LRC suportados, retirando os marcadores de tempo sem reordenar ou remover repetições. A aceitação considerou conflitos entre textos simples e sincronizados, bem como conflitos entre candidatos elegíveis. Revisões dirigidas examinaram diferenças de caixa, variantes documentadas de artista e consultas inicialmente interrompidas pelo limite de candidatos. Os vínculos preservam o texto original e sua proveniência por identificador e hash. A verificação de integridade confirma o conteúdo armazenado e utilizado. A correção da correspondência depende das regras de identidade e das evidências disponíveis para cada caso.

\begin{table}[H]\centering\small
\caption{Desfechos da vinculação no sorteio original de 9.526 faixas.}\label{tab:status_nova}
\begin{tabular}{@{}P{.22\linewidth}rP{.60\linewidth}@{}}\toprule
Desfecho & Faixas & Critério operacional \\\midrule
Letra aceita & 3.869 & Vínculo elegível com texto utilizável, incluindo reutilizações compatíveis. \\
Ambígua & 1.085 & Evidência insuficiente para resolver candidatos, versões ou textos concorrentes. \\
Sem candidato & 4.155 & Ausência de candidato pelos métodos e fontes consultados. \\
Instrumental & 416 & Indicação no registro do LRCLIB, distinta do atributo acústico do Spotify. \\
Sem texto & 1 & Registro compatível, sem letra utilizável. \\\midrule
Total & 9.526 & Sorteio original do catálogo. \\\bottomrule
\end{tabular}
\end{table}

Houve 3.869 vínculos aceitos, com cobertura de 40,62\% sem pesos e 40,63\% com os pesos do sorteio (Tabela~\ref{tab:status_nova}). A cobertura ponderada é uma estimativa da fração alcançada pelo procedimento de recuperação nas condições observadas. Ela não mede toda letra existente. A Tabela~\ref{tab:cobertura_nova} mostra o comportamento por banda de popularidade. Nos gêneros \textit{sleep} e \textit{detroit-techno}, nenhuma das 147 e 32 faixas sorteadas teve letra aceita, respectivamente. Em \textit{indie-pop}, foram 63 de 82. Essas diferenças explicam por que os vínculos aceitos mudam a composição da amostra. Em particular, a popularidade média passou de 18,38 no catálogo para 24,11 nas 3.869 aceitações e a instrumentalidade média, de 0,252 para 0,107.

\begin{table}[H]\centering\small
\caption{Cobertura no sorteio original por banda de popularidade.}\label{tab:cobertura_nova}
\begin{tabular}{lrrr}\toprule
Escore & Sorteadas & Aceitas & Cobertura (\%) \\\midrule
0--19 & 5.548 & 1.664 & 29,99 \\
20--39 & 2.919 & 1.565 & 53,61 \\
40--59 & 935 & 568 & 60,75 \\
60--79 & 119 & 71 & 59,66 \\
80--100 & 5 & 1 & 20,00 \\\midrule
Total & 9.526 & 3.869 & 40,62 \\\bottomrule
\end{tabular}
\nota{Cobertura calculada dentro de cada banda. A última banda tem cinco faixas.}
\end{table}

Um piloto online em 100 faixas sem candidato retornou zero vínculos textuais provisórios sob os critérios usados. O desfecho caracteriza o rendimento daquele piloto. A recuperação seletiva motivou construir a coorte de modelagem a partir de um quadro mais amplo de vínculos aceitos.

\subsection{Do quadro vinculado às 9.526 faixas classificadas}
A Tabela~\ref{tab:trajeto} distingue as etapas. O sorteio inicial continua sendo a referência para dimensionamento e cobertura. Para o modelo, reuniu-se um quadro de 80.943 faixas com letra, combinando as aceitações do sorteio atual com registros de uma etapa anterior de recuperação e com revisões dirigidas. A etapa anterior forneceu 77.535 vínculos. Dos 3.869 do sorteio atual, 682 já estavam nesse conjunto. Somaram-se 42 vínculos de revalidações históricas e 179 de uma busca por gêneros deficitários. A procedência anterior importa para explicar a formação do quadro e a exposição de faixas em análises anteriores.

\begin{table}[H]\centering\small
\caption{Percurso das faixas desde o catálogo até a avaliação dos modelos.}\label{tab:trajeto}
\begin{tabular}{@{}P{.36\linewidth}rP{.49\linewidth}@{}}\toprule
Conjunto & Faixas & Papel no estudo \\\midrule
Catálogo Spotify & 1.159.764 & População finita descrita e origem do sorteio. \\
Sorteio inicial & 9.526 & Seleção probabilística para consultar letras e medir cobertura. \\
Vínculos no sorteio inicial & 3.869 & Parcela do sorteio com letra aceita. \\
Quadro consolidado com letra & 80.943 & Vínculos de recuperação atual, anterior e revisões dirigidas. \\
Coorte de modelagem & \textbf{9.526} & \textbf{Faixas com letra avaliadas nos quatro modelos.} \\\bottomrule
\end{tabular}
\end{table}

No quadro de vínculos, a alocação foi definida nas células de gênero e das cinco bandas da Tabela~\ref{tab:cobertura_nova}. Uma programação inteira manteve os totais das bandas do sorteio original, garantiu ao menos uma faixa em cada célula com capacidade e preservou a cota de gênero até o limite viável. A redistribuição priorizou reduzir os desvios absolutos nas cotas de gênero e, em seguida, nas células. Após fixar as cotas, sorteou-se sem reposição em cada célula, com semente 20261108. Para uma célula $h$ com capacidade $N_h^L$ e alocação $n_h^L$, a probabilidade de inclusão condicional ao quadro é $n_h^L/N_h^L$. Os pesos $w_i^L=N_h^L/n_h^L$ são usados em uma análise de sensibilidade da avaliação.

A coorte tem 9.526 IDs distintos e 82 gêneros. A conferência dos IDs mostrou que 461 faixas da coorte vieram do sorteio inicial e 9.065 vieram de fora dele. A coorte preserva o tamanho e os totais das cinco bandas do sorteio original. Restaram diferenças de gênero, especialmente em \textit{sleep} (23 faixas, ante cota de 147) e \textit{indie-pop} (81, ante 82). O escore médio da coorte é 19,87. Seu valor pode mudar sem alterar os totais por banda, pois há variação dentro das bandas. A instrumentalidade média é 0,135 e a energia média, 0,673, ante 0,252 e 0,640 no catálogo. A Figura~\ref{fig:selecao_final} apresenta os desvios etapa a etapa.

\figura[.37\textheight]{selecao_metadados_final}{Desvios das médias em relação ao catálogo nas três etapas de seleção.}{fig:selecao_final}{Sorteio inicial ponderado e demais etapas com peso igual por faixa. A unidade é o desvio-padrão do catálogo.}

\begin{table}[H]\centering\small
\caption{Formatos presentes nos registros de origem das letras usadas na modelagem.}\label{tab:formatos_coorte}
\begin{tabular}{lrr}\toprule
Formato registrado & Faixas & Percentual (\%) \\\midrule
Somente simples & 3.726 & 39,11 \\
Somente sincronizada & 0 & 0,00 \\
Simples e sincronizada & 5.800 & 60,89 \\
Total & 9.526 & 100,00 \\
\bottomrule\end{tabular}\nota{Denominador por faixa Spotify.}
\end{table}


Nos registros associados à coorte, 5.800 das 9.526 faixas têm as duas formas de letra registradas e 3.726 têm somente a forma simples (Tabela~\ref{tab:formatos_coorte}). A tabulação é por faixa Spotify, vinculada ao registro de origem. Não equivale a uma contagem de letras distintas nem a uma estimativa de formatos no dump inteiro. A formação da coorte permite comparar áudio e letras nas mesmas observações. O alcance fica condicionado ao quadro recuperado, porque os pesos do segundo sorteio descrevem probabilidades de inclusão naquele quadro.

